% TOP_PROBLEM: {"problemId": "problem.grothendieck-standard-conjecture-d", "canonicalTitle": "Grothendieck's Standard Conjecture D", "releaseRank": 44, "primaryDomain": "number_theory_arithmetic_geometry", "primaryDomainLabel": "Number theory & arithmetic geometry", "displayStatus": "open_with_solved_subcases", "releaseStatus": "open_with_solved_subcases"}
% !TeX program = lualatex
% Definition and sources only; this file is not an LLM solution attempt.
\documentclass[11pt]{article}
\usepackage[margin=25mm]{geometry}
\usepackage{fontspec}
\setmainfont{Latin Modern Roman}
\usepackage{amsmath,amssymb,mathtools}
\usepackage{unicode-math}
\setmathfont{Latin Modern Math}
\usepackage{xurl,hyperref}
\hypersetup{colorlinks=true,urlcolor=blue}
\setcounter{secnumdepth}{0}
\setlength{\parindent}{0pt}
\setlength{\parskip}{5pt}
\setlength{\emergencystretch}{3em}
\begin{document}
\section*{\texorpdfstring{Grothendieck's Standard Conjecture D}{Grothendieck's Standard Conjecture D}}\label{44}
\subsection{Definitions and mathematical statement}
A Weil cohomology theory is a finite-dimensional graded cohomology theory on smooth projective varieties with cycle classes, products, Poincaré duality, and the standard compatibility axioms. A rational cycle $z$ of codimension $r$ on an $n$-fold is numerically trivial if $\deg(z\cdot w)=0$ for every complementary cycle $w$. It is homologically trivial if $\operatorname{cl}_H(z)=0$. The assertion is
\[
 \bigl[\forall w\in\operatorname{CH}^{n-r}(X)_{{\mathbb{Q}}},\ \deg(z\cdot w)=0\bigr]
 \quad\Longrightarrow\quad \operatorname{cl}_H(z)=0
\]
for every field, available Weil cohomology theory, smooth projective variety, and codimension in the recorded scope.
\par\smallskip\textit{Formulation and concept references:} \href{https://arxiv.org/pdf/1905.04626}{[S1]}.

\subsection{Short English statement}
If an algebraic cycle has zero intersection number with every complementary cycle, must its cohomology class vanish?
\subsection{Sources}
\begin{itemize}
\item[S1]Michael K. Brown and Mark E. Walker, Standard Conjecture D for matrix factorizations, classical Conjecture1.1 and surrounding introduction.
\url{https://arxiv.org/pdf/1905.04626}.
\textit{Formulation source.}
\item[S2]The Tate conjecture for abelian fourfolds over finite fields.
\url{https://arxiv.org/abs/2608.28651}.
\textit{Further reference; not a proof endorsement.}
\item[S3]Tate's question, Standard conjecture D, semisimplicity and Dynamical degree comparison conjecture.
\url{https://arxiv.org/abs/2503.09432}.
\textit{Further reference; not a proof endorsement.}
\end{itemize}
\end{document}
